The transition toward sustainable and low-carbon energy systems has driven the widespread adoption of distributed combined cooling, heating and power (CCHP) systems, which achieve higher primary energy utilization through the cascade use of multiple energy forms \cite{liu2018combined,cho2009combined}. Unlike conventional separate production systems, CCHP systems integrate prime movers such as gas turbines with thermally activated equipment including absorption chillers and heat recovery units, enabling simultaneous provision of electricity, heating, and cooling from a single fuel source \cite{wu2016multiobjective}. The planning of such systems---determining the optimal capacity of each component---is a critical decision that fundamentally shapes both the economic performance and operational flexibility over the system's lifetime \cite{li2016optimal}.

The optimal planning of CCHP systems has been extensively studied within the broader framework of integrated energy systems (IES) \cite{mancarella2014mes,geidl2007energy}. Early studies predominantly adopted single-objective optimization approaches, minimizing either the total annualized cost or the primary energy consumption \cite{ren2010optimal}. Recognizing the multi-dimensional nature of energy system performance, subsequent research has formulated multi-objective optimization models that simultaneously consider economic, environmental, and reliability criteria \cite{wang2019multi,wang2015modelling}. The non-dominated sorting genetic algorithm II (NSGA-II) has become a widely adopted solver for such problems due to its effectiveness in generating well-distributed Pareto fronts \cite{deb2002fast}. However, the majority of existing multi-objective formulations pair economic cost with carbon emissions or energy efficiency, while the temporal coordination between energy supply and demand---referred to as source-load matching---has received comparatively limited attention as an explicit optimization objective.

Source-load matching characterizes the degree to which the temporal profile of distributed energy supply aligns with the corresponding demand profile. Poor matching leads to excessive grid interaction, curtailment of renewable generation, and underutilization of installed capacity, all of which undermine the intended benefits of distributed energy systems. Several indices have been proposed to quantify this matching performance. The net load volatility, measured by the standard deviation of the difference between supply and demand, captures the temporal fluctuation of the supply-demand imbalance \cite{zhang2020comprehensive}. The Pearson correlation coefficient evaluates the linear co-movement between supply and demand time series but is insensitive to magnitude differences. The self-sufficiency rate measures the fraction of demand met by local generation but ignores temporal distribution. While these indices provide useful perspectives, they share a common limitation: they treat electricity, heating, and cooling as equivalent energy forms, assigning equal weight to mismatches across all three carriers.

From a thermodynamic standpoint, different energy forms possess fundamentally different qualities, or abilities to perform useful work. Electricity, as pure exergy, represents the highest quality energy carrier, while low-temperature heating and cooling have significantly lower exergy content---typically only 10--15\% of the work potential of an equivalent amount of electrical energy \cite{dincer2001energy,bejan2016advanced,hu2020eqc}. The concept of energy quality coefficient (EQC), defined as the ratio of exergy to energy, provides a rigorous framework for quantifying these differences \cite{chen2022energy,ma2023exergyflex}. In the context of IES optimization, exergy-based analysis has gained increasing attention. Li et al.\ established a unified exergy flow calculation model for regional IES, demonstrating that the spatial distributions of energy quantity and energy quality exhibit fundamentally different patterns \cite{li2022exergyflow}. Chen et al.\ proposed a joint solution framework for IES optimal operation and evaluation based on exergy, revealing that energy-based pricing underestimates the value of electrical loads by up to 47.56\% compared to exergy-based pricing \cite{chen2024joint}. Ma et al.\ introduced an ``equal exergy replacement'' mechanism for distributed multi-energy systems, establishing the principle that high-quality energy must correspond to high economic value \cite{ma2023exergyflex}. More recently, Duo et al.\ proposed a multi-objective IES planning method that simultaneously considers energy quantity and quality matching, achieving 2.80\% higher efficiency and 10.85\% lower cost \cite{duo2025matching}. These studies have established the value of exergy analysis in IES operational optimization and evaluation. However, the application of exergy analysis to \emph{capacity planning}---specifically, using Carnot exergy factors as energy quality weights in source-load matching quantification---remains unexplored. Ignoring energy quality differences in matching assessment can lead to suboptimal planning decisions: equal weighting implicitly treats 1~kW of electrical mismatch as equivalent to 1~kW of thermal mismatch---a premise that contradicts the second law of thermodynamics \cite{moran2014fundamentals}---and may cause the optimizer to over-invest in low-grade thermal storage when the same investment in electrical storage would yield greater overall system benefit.

Meanwhile, the emergence of Carnot battery technology offers new possibilities for enhancing source-load matching in CCHP systems. A Carnot battery converts surplus electricity into thermal energy via a heat pump during charging and recovers electricity from stored heat via a heat engine during discharging \cite{dumont2022carnot,novotny2022review}. Unlike conventional single-carrier storage devices, the Carnot battery inherently couples electrical and thermal energy domains. Recent studies have demonstrated that Carnot battery-based multi-energy systems with cascaded latent thermal stores can achieve comprehensive COP up to 95\% and exergy efficiency of 55\% in CCHP mode, simultaneously delivering multi-grade heat and cold alongside electricity \cite{zhao2024cbthermo,huang2025cbdynamic}. A comprehensive review confirms that multi-output Carnot batteries exhibit higher energy and exergy efficiencies as well as improved economics compared to single-output configurations \cite{lykas2026cbreview}. Furthermore, thermally integrated Carnot batteries combined with district heating have been shown to reduce heating station capacity by up to 47\% \cite{poletto2024cbdh}, and techno-economic analyses indicate that the electro-thermal coupling structure is cost-optimal for residential PV-heat pump systems \cite{laterre2025cbresidential}. Despite these advances, existing research has predominantly focused on component-level thermodynamic optimization or dynamic control strategies. The integration of Carnot batteries into CCHP capacity planning---particularly in conjunction with source-load matching optimization---has not been investigated.

To address these research gaps, this paper makes two mutually reinforcing contributions. First, an energy-quality-weighted three-dimensional Euclidean distance index is proposed for quantifying source-load matching in CCHP systems. Unlike previous studies that adopt empirically adjusted energy quality coefficients, and distinct from existing exergy-based approaches that focus on operational optimization and evaluation \cite{li2022exergyflow,chen2024joint,ma2023exergyflex}, this work applies Carnot exergy factors to a fundamentally different problem: weighting the relative importance of different energy carriers in a planning-stage matching index. The resulting coefficients are parameter-free, uniquely determined by supply temperatures and ambient conditions of each case (e.g., $\lambda_e = 1.0$, $\lambda_h = 0.131$, $\lambda_c = 0.064$ for the German case). Post-optimization operational analysis confirms that the Carnot-factor-weighted index produces solutions with superior grid-interaction characteristics---including 26.7\% lower peak grid demand and 47.5\% higher self-sufficiency rate compared to the conventional volatility-based method at equivalent budget levels. Second, a Carnot battery model with waste heat recovery is integrated into the CCHP system topology as a representative electro-thermal coupled storage device, and its capacity is co-optimized with conventional equipment within the bi-objective framework. The Carnot battery's inherent cross-carrier coupling---converting electricity to heat during charging and recovering electricity from heat during discharging \cite{zhao2024cbthermo}---serves as a natural testbed for evaluating the proposed energy-quality-weighted index: only an index that correctly differentiates the thermodynamic quality of electrical and thermal energy can properly guide the sizing of such coupled devices. The two contributions thus form a closed validation loop: the theoretical innovation (exergy-weighted matching) enables effective integration of the technological innovation (Carnot battery), and the technological innovation reciprocally validates the theoretical foundation. The optimization employs NSGA-II with a nested operational dispatch model solved by the Open Energy Modelling Framework (oemof) \cite{hilpert2018oemof}, using a typical-day decomposition approach to maintain computational tractability \cite{kotzur2018impact,dominguez2012selection}.

The proposed framework is validated through numerical experiments across two representative cases: a heating-dominated German community and a cooling-dominated campus in Songshan Lake, China. Five matching indices are systematically compared, and the practical value of the proposed index is validated through full 8760-hour post-optimization operational analysis. The remainder of this paper is organized as follows. \Cref{sec:system} presents the CCHP system topology and component models. \Cref{sec:matching} develops the energy-quality-weighted matching index with theoretical justification. \Cref{sec:optimization} formulates the bi-objective optimization framework. \Cref{sec:cases} describes the two case studies, and \Cref{sec:results} presents and discusses the results. \Cref{sec:conclusions} summarizes the key findings and outlines future research directions.
